\section{Privacy and Threat-Model Enforcement}
\label{sec:impl-privacy-enforcement}

The task description (Page~iv) names a strict zero-egress policy and a
threat model that includes retrieval poisoning. This section describes the
three enforcement layers in the delivered system: a network-policy module
that gates every outbound socket, a signed-corpus verifier that detects
RAG knowledge-base tampering at load time, and a containerised harness
that pins the resource and network namespace.

\subsection{Network Policy and Zero-Egress Gate}
\label{sec:impl-network-policy}

The extension routes every outbound HTTP(S) call through a single
chokepoint, \texttt{src/networkPolicy.ts}. The module exports an
\texttt{assertEgressAllowed(url)} function that throws an
\texttt{EgressBlockedError} unless the URL's hostname is on a fixed
loopback allowlist (\texttt{127.0.0.1}, \texttt{::1}, \texttt{localhost},
\texttt{0.0.0.0}) \emph{or} the user has explicitly opted in via the new
\texttt{codeGuardian.enableExternalDataFetch} setting (default
\texttt{false}). The opt-in path exists for the public CWE/CVE/OWASP
metadata refresh; it does not relax the policy for any source-code-bearing
call. \texttt{src/vulnerabilityDataManager.ts} consults the gate before
every fetch and additionally treats any cached file as valid when the gate
is closed, so a fully offline machine can still load the bundled
vulnerability snapshot. Because the gate is a synchronous assertion at
the call site, it cannot be bypassed by code that holds a live HTTP
client reference.

\subsection{Signed Corpus and Provenance}
\label{sec:impl-corpus-signing}

Retrieval poisoning is the threat that an attacker swaps the on-disk RAG
knowledge base for one carrying malicious ``guidance'' that the LLM then
trusts as authoritative context. To detect this attack, the RAG corpus is
signed with an Ed25519 detached signature over a SHA-256 manifest of every
corpus file. The signing tool
(\texttt{evaluation/sign-corpus.js}) generates the keypair on first run,
walks \texttt{security-knowledge/}, computes a SHA-256 for each file,
attaches per-file provenance from
\texttt{evaluation/corpus-provenance.json} (source URL plus retrieval
timestamp), produces a canonical tab-separated payload sorted by file
path, signs the payload, and writes
\texttt{security-knowledge/corpus-manifest.json}. The private key never
ships with the extension; the public key
(\texttt{keys/corpus-public.pem}) ships in the package and is bundled
into the Docker image.

At RAG load time \texttt{src/corpusVerifier.ts} re-derives the manifest
hashes from the on-disk content, recomputes the canonical payload, and
verifies the Ed25519 signature against the public key. A hash mismatch,
a missing file, or a signature failure makes the load abort under
the default \texttt{codeGuardian.requireSignedCorpus=true} setting.
Because provenance fields are bound into the signed payload, an attacker
cannot rewrite a corpus file's ``retrieved from'' string without
invalidating the signature.

\subsection{Containerised Harness}
\label{sec:impl-container}

\texttt{Dockerfile} pins \texttt{node:20.19.0-alpine}, installs the
locked dependency set with \texttt{npm ci}, copies the harness, the signed
corpus, and the public key into the image, and compiles the extension at
build time. The private key is excluded by \texttt{.dockerignore} so
images cannot accidentally carry signing material.
\texttt{docker-compose.yml} brings up Ollama and the harness on a
user-defined bridge network with explicit CPU and memory limits (8 CPUs,
16\,GB) and binds the Ollama port to loopback only. This is the
``containerised execution'' clause of the task description's
reproducibility triad and the network-isolation arm of the threat model
in one stack: the harness has no path to a non-loopback host, and resource
measurements (Section~\ref{sec:eval-r4-resources}) are comparable across
hosts because the container caps the harness side identically.
